% Inserted after Cor: the floor is not attained in finite time.
% Relay drift: a corollary of the composition law, in a transmission setting.

\subsection{Relayed content drifts multiplicatively}
\label{sec:relay}

The composition law was stated for a cascade of events acting on one
observer. The same algebra governs a second situation, which we record
because the setting differs enough that the identity is easy to miss: an
account passed from one observer to another along a chain.

\begin{definition}[Relay chain]
\label{def:relay}
A \emph{relay} is an observer that receives an account, individuates it
within its own graph, and emits the result. A \emph{relay chain} of length
$n$ is a sequence of relays $\gamma_1,\dots,\gamma_n$, each applying its own
individuation with power $\kap_i$ to the account it received.
\end{definition}

\begin{corollary}[Relayed content drifts multiplicatively]
\label{cor:relay-drift}
Along a relay chain the surviving fidelity of an account is
\[
\prod_{i=1}^{n}\bigl(1-\kap(\gamma_i)\bigr),
\]
the residual factor of \Cref{thm:composition}, and the accumulated drift is
$1-\prod_i(1-\kap_i)$. If the relay powers are bounded below by $c>0$ the
fidelity decays geometrically, at most $(1-c)^{n}$ after $n$ relays.
\end{corollary}

\begin{proof}
Each relay individuates the account it received, so by \Cref{def:power} its
power is the fraction of the outstanding gap it closes, and by
\Cref{def:resfactor} the fraction surviving it is $1-\kap_i$. The account
reaching relay $i+1$ is the account surviving relay $i$, so the surviving
fractions compose exactly as in the proof of \Cref{thm:composition}, giving
the stated product. The geometric bound is \Cref{prop:diminishing}.
\end{proof}

\begin{remark}[What the corollary does and does not say]
\label{rem:relay-scope}
Two readings must be separated, because only one is licensed.

The corollary says that drift accumulates \emph{without any relay altering
the account on purpose}. Each relay performs an ordinary individuation; the
drift is the composition of those individuations and requires no
intervention, no intent, and no defect in any relay. A chain of faithful
relays produces a drifted account.

The corollary does not say that the terminal account is false, or that any
relay erred. By \Cref{cor:nosharp} no account resolves to a point in any
observer's graph, so there is no value against which a relayed account could
be compared to establish falsity. Whether two accounts stand in some
relation to each other is a question an outside party may compute from both;
it is not a property either relay holds, and nothing in this algebra assigns
one.

The distinction matters for instrumentation. A system that records relay
depth can report accumulated drift, which is a computable quantity; a system
that reports an account as corrupted is asserting a comparison it has not
performed.
\end{remark}

\begin{remark}[Provenance depth is worth recording]
\label{rem:relay-depth}
\Cref{cor:relay-drift} gives a reason to carry relay depth alongside a
value. Depth alone bounds the accumulated drift when a lower bound on relay
powers is known, and does so without inspecting the content: a value at
depth $n$ has fidelity at most $(1-c)^{n}$. This is a cheap invariant and it
is the only statement about a relayed account that the algebra supports
without further data.
\end{remark}
